\centerline{\bf \Huge Алгебра для самооценки}

\begin{flushright}
        \scriptsize{Единственный счастливый предмет, который вам понадобится, уже находится внутри — УВЕРЕННОСТЬ\\ Мидорима Синтаро}
\end{flushright}

\problem{1}
Существуют ли 2024 целых чисел, сумма и произведение которых равны 2024?

\problem{2}
Существуют ли нецелые числа $x$ и $y$, для которых $\{x\}\{y\}=\{x+y\}$ ?

\problem{3}
Выясните, рационально ли число

$$
\alpha=\sqrt[3]{4+\sqrt{15}}+\sqrt[3]{4-\sqrt{15}}
$$
%Прасолов 6.18

\problem{4} Найдите все решения системы уравнений:

$\left\{\begin{array}{l}
    x(y+z)=35 \\ 
    y(x+z)=32 \\ 
    z(x+y)=27
    \end{array}\right.$
%Прасолов 3.1

\problem{5}
Десять рабочих должны собрать из деталей 50 изделий. Сначала детали каждого изделия нужно покрасить; на это одному рабочему требуется 10 минут. После окраски детали 5 минут сохнут. На сборку изделия одному рабочему требуется 20 минут. Сколько рабочих нужно назначить малярами и сколько монтажниками, чтобы работа была выполнена в кратчайшее время? (Двое рабочих не могут одновременно красить или монтировать одно и то же изделие.)

%Прасолов 7.16

\problem{6}
Назовём \textit{паспортом} натурального числа набор простых чисел, входящих в его разложение. Докажите, что в любой натуральной арифметической прогресии есть сколько угодно много чисел с одинаковым паспортом.

\problem{7}
Последовательности $\left\{a_n\right\},\left\{b_n\right\}$ таковы, что для каждого натурального $n$

$$
a_n>0, \quad b_n>0
$$
и
$$
a_{n+1}=a_n+\frac{1}{b_n}, \quad b_{n+1}=b_n+\frac{1}{a_n}
$$


Докажите, что $a_{50}+b_{50}>20$.

%https://math.mosolymp.ru/upload/files/2025/khamovniki/9/2024-09-16_Ochnyiy_otbor,_resheniya.pdf

\problem{8$^{*}$}
Докажите, что если $r$ это простой делитель числа $\frac{x^p-1}{x-1}$, где $p$ простое число, то либо $r \equiv 1\ (mod\  p)$ либо $r=p$.